\section{Christ, Mary, and the Angels}\label{sec:christ-mary}

The angels belong to Christ. They were made through him and for him; they
adore him when the Father brings the First-begotten into the world; they
announce his conception, sing at his birth, serve him in the desert,
strengthen him in his agony, keep his tomb, escort his Ascension, and will
come with him when he returns to judge. The Apostle names their orders in
order to set the Son above them: ``in him were all things created in heaven
and on earth, visible and invisible, whether thrones, or dominations, or
principalities, or powers. All things were created by him and in him'' (Col
1:16). The Catechism gathers the whole into one sentence, ``Christ is the
centre of the angelic world,'' and gives the reason: the angels are his
because they were created through and for him, and still more because he
has made them messengers of his saving plan (CCC~331). The Mother of this
Lord, saluted by an angel and served by angels, is exalted above them all;
and the men whom Christ redeems are taken up into their company. The
Fathers pass through each of these mysteries with the angels at their side,
and Aquinas gives three of them their own articles: Christ's headship over
the angels (\ST{III q.~8 a.~4}), his judicial power over them (\ST{III
q.~59 a.~6}), and the Annunciation (\ST{III q.~30}).

\sectionguard
\subsection{Christ, Head of the Angels}\label{sec:cm-head}

\paragraph{The witness of the Apostle.} Twice the Apostle sets Christ over
the orders by name. To the Colossians he says that Christ is he ``who is
the head of all principality and power'' (Col 2:10). To the Ephesians he
describes the power the Father ``wrought in Christ, raising him up from the
dead and setting him on his right hand in the heavenly places. Above all
principality and power and virtue and dominion and every name that is
named, not only in this world, but also in that which is to come. And he
hath subjected all things under his feet and hath made him head over all
the church'' (Eph 1:20--22). Peter says the same of the ascended Christ:
``being gone into heaven, the angels and powers and virtues being made
subject to him'' (1~Pet 3:22). The Father's design is ``to re-establish all
things in Christ, that are in heaven and on earth, in him'' (Eph 1:10), and
the name of Jesus is the name at which ``every knee should bow, of those
that are in heaven, on earth, and under the earth'' (Phil 2:10).

The Epistle to the Hebrews opens its doctrine of the Son with a comparison
to the angels. Christ is ``made so much better than the angels as he hath
inherited a more excellent name than they''; to none of them did the Father
say, ``Thou art my Son''; and ``when he bringeth in the first begotten into
the world, he saith: And let all the angels of God adore him'' (Heb 1:4--6;
Ps 96[97]:7). The angels are made ``spirits'' and ``a flame of fire,'' and
they are ``all ministering spirits, sent to minister for them who shall
receive the inheritance of salvation'' (Heb 1:7, 14; Ps 103[104]:4). The
second chapter gives the other side of the mystery. Jesus ``was made a
little lower than the angels, for the suffering of death, crowned with glory
and honour,'' and ``nowhere doth he take hold of the angels: but of the seed
of Abraham he taketh hold'' (Heb 2:9, 16). ``God hath not subjected unto
angels the world to come'' (2:5). The Son is above the angels by nature and
by inheritance, beneath them for a time in the flesh he assumed for death,
and set over them in that same flesh, crowned after the Passion.

\paragraph{The Fathers on the Head above the orders.} Irenaeus reads the
headship as the completion of the Word's work of summing up. The Word made
man took up man into himself, ``thus summing up all things in Himself: so
that as in super-celestial, spiritual, and invisible things, the Word of God
is supreme, so also in things visible and corporeal He might possess the
supremacy, and, taking to Himself the pre-eminence, as well as constituting
Himself Head of the Church, He might draw all things to Himself at the
proper time'' (\emph{Adv.\ haer.}\ III.16.6). The Word was already supreme
over the invisible creation; the Incarnation extends the same primacy
through the visible creation and gathers both under one Head.

Chrysostom is the preacher of this exaltation. On the command that the
angels adore the First-begotten, he observes that Paul speaks ``concerning
that which is according to the flesh'': the Father, ``just as any one
introducing another into a house straightway commands those having the care
thereof to do him reverence,'' says ``in regard to the Flesh, `And let all
the Angels of God worship Him'\,'' (\emph{Hom.\ in Heb.}\ 3.1). On Heb 2:16
he reads the verb ``taketh hold'' as the figure of a pursuer overtaking one
who flees: human nature was running from God, and he ran after it and
caught it. Then he cries out that it is ``a great and a wonderful thing, and
full of amazement that our flesh should sit on high, and be adored by Angels
and Archangels, by the Cherubim and the Seraphim'' (\emph{Hom.\ in Heb.}\
5.1; \S\ref{sec:fb-chrysostom}). On Ephesians he dwells on the Apostle's ``far above.'' ``He does
not simply say, `above,' but, `far above;' for God is above those powers
which are above. And thither then hath He raised Him, Him that is one of
us, brought Him from the lowest point to the supremest sovereignty, to that
beyond which there is no other honor'' (\emph{Hom.\ in Eph.}\ 3). The
exaltation of the Head is a charge laid on the members: ``Yet Angels
reverence that Head, and Archangels, and all those powers above. And shall
we, which are His body, be awed neither on the one account nor the other?''
(\emph{ibid.}). The body whose Head the angels reverence cannot be handed
over to the insults of the devils.

\paragraph{The angels in the mystical body (\ST{III q.~8 a.~4}).} In the
Tertia Pars Aquinas treats the grace of Christ as Head of the Church, and
asks whether Christ, as man, is also the Head of the angels. Three
objections stand against it. Head and members must be of one nature, and
Christ as man shares his nature with men alone, for ``nowhere doth He take
hold of the angels'' (Heb 2:16). Christ is Head of the Church, and the
Church is the congregation of the faithful; but the angels have no faith,
for they walk ``by sight'' and not ``by faith'' (2~Cor 5:6--7). And
Augustine says that as the Word quickens souls, the Word made flesh quickens
bodies, which the angels lack. The \emph{sed contra} is the Apostle's
``head of all Principality and Power'' (Col 2:10), ``and the same reason
holds good with the other orders of angels.''

Aquinas answers from the nature of a body. Where there is one body there is
one head, and ``a multitude ordained to one end, with distinct acts and
duties, may be metaphorically called one body. But it is manifest that both
men and angels are ordained to one end, which is the glory of the Divine
fruition. Hence the mystical body of the Church consists not only of men
but of angels'' (\latin{corpus Ecclesiae mysticum non solum consistit ex
hominibus, sed etiam ex Angelis}). ``Now of all this multitude Christ is the
Head, since He is nearer God, and shares His gifts more fully, not only than
man, but even than angels; and of His influence not only men but even
angels partake.'' The proof is Eph 1:20--22, and the sign is the Gospel's
``angels came and ministered to Him'' (Matt 4:11). The replies carry the
doctrine further. Christ's influence reaches men chiefly in their souls,
``wherein men agree with angels in generic nature, though not in specific
nature'' (ad~1). ``The Church, on earth, is the congregation of the
faithful; but, in heaven, it is the congregation of comprehensors. Now
Christ was not merely a wayfarer, but a comprehensor. And therefore He is
the Head not merely of the faithful, but of comprehensors'' (ad~2). And the
humanity of Christ, by virtue of the Godhead to which it is joined, ``can
cause something not only in the spirits of men, but also in the spirits of
angels, on account of its most close conjunction with God, i.e.\ by
personal union'' (ad~3). The Church of which Christ is Head is one city of
men and angels, and his grace flows into both.

\paragraph{The Judge of angels (\ST{III q.~59 a.~6}).} The last article of
the treatise on the mysteries of Christ asks whether his judicial power
extends to the angels. The objections: the angels were judged once for all
at the beginning, when some fell and the rest were confirmed; the angels
will come to judge with Christ, ``when the Son of Man shall come in His
majesty, and all the angels with Him'' (Matt 25:31), and one cannot be both
judge and judged; and to judge the angels would be to judge all creatures,
which belongs to God's providence. The \emph{sed contra} is Paul's question,
``Know you not that we shall judge angels?'' (1~Cor 6:3); ``the saints
judge only by Christ's authority,'' and much more, then, does Christ judge
them.

Aquinas determines that ``the angels are subjects of Christ's judiciary
power, not only with regard to His Divine Nature, as He is the Word of God,
but also with regard to His human nature,'' and gives three reasons. The
first is the closeness of the assumed nature to God: Christ's soul ``is more
filled with the truth of the Word of God than any angel: for which reason He
also enlightens the angels, as Dionysius says (Coel.\ Hier.\ vii), and so He
has power to judge them.'' The second is the merit of the Passion: ``by the
lowliness of His Passion, human nature in Christ merited to be exalted above
the angels,'' so that every knee bows at the name of Jesus (Phil 2:10), and
``in token whereof it is said in the Apocalypse (7:11) that `all the angels
stood round about the throne.'\,'' The third is what the angels do for men,
of whom Christ is Head in a special way: they are ``ministering spirits''
sent for the heirs of salvation (Heb 1:14). The angels are subject to his
judgment in three respects: in the dispensing of what is done through
them, which the Man Christ also dispenses, to whom angels ministered and
from whom the demons begged to be sent into the swine; in their accidental
rewards, such as ``the joy which they have at the salvation of men,
according to Luke 15:10,'' and in the accidental punishments of the demons,
as when the demon cried, ``What have we to do with thee, Jesus of Nazareth?
art Thou come to destroy us?'' (Mark 1:24); and in the essential reward of
beatitude and the essential punishment of damnation, which ``was done by
Christ from the beginning of the world, inasmuch as He is the Word of God.''
The replies answer the objections in turn. The judgment at the beginning
concerned the essential reward (ad~1). The angels judge as spiritual
creatures, yet, as Augustine says, ``Although the spiritual man judgeth all
things, still he is judged by Truth Itself'' (ad~2; \emph{De vera rel.}\
31). And since the lower things are ruled by God through the higher, ``it
must be said that all things are ruled by Christ's soul, which is above
every creature'' (ad~3); for this reason ``God hath not subjected unto
angels the world to come'' (Heb 2:5).

\paragraph{The angels and the mystery of the Incarnation.} The Apostle says
that the great mystery of godliness ``appeared unto angels'' (1~Tim 3:16),
that the manifold wisdom of God is ``made known to the principalities and
powers in heavenly places through the church'' (Eph 3:10), and Peter calls
the Gospel a thing ``on whom the angels desire to look'' (1~Pet 1:12).
Aquinas explains how the angels know it (\ST{I q.~57 a.~5}). By their
natural knowledge they cannot know the mysteries of grace, since these
``depend upon the pure will of God.'' In the vision of the Word they know
them, ``but not all mysteries: nor do they all know them equally,'' and ``the
higher angels beholding the Divine wisdom more clearly, learn more and
deeper mysteries in the vision of God, which mysteries they communicate to
the lower angels by enlightening them. Some of these mysteries they knew
from the very beginning of their creation; others they are taught
afterwards, as befits their ministrations.'' The Incarnation in general was
revealed to all of them from the beginning of their beatitude, ``a kind of
general principle to which all their duties are ordered'' (ad~1); its
special conditions even the highest learned later. The \emph{sed contra} is
Dionysius's reading of Isaias 63, where the heavenly essences ask, ``Who is
this that cometh from Edom?'' and are taught by Jesus himself. Dionysius
marvels that ``even the first of the Beings in Heaven, and so far above
all, should reverently strive after the supremely Divine illuminations'':
they ask one another first, and then Jesus teaches them ``immediately, and
shewing to them, at first hand, His beneficent work out of love to man''
(\emph{CH}~7.3; see \S\ref{sec:ch7}). The fallen angels knew far less. The
demons ``much less fully understood the mystery of the Incarnation, when
Christ was in the world,'' for, as Augustine says, ``It was not manifested
to them as it was to the holy angels~\dots\ but it was made known by some
temporal effects, so as to strike terror into them''; had they known,
``they would never have procured the crucifixion of the Lord of glory''
(\ST{I q.~64 a.~1 ad~4}; 1~Cor 2:8; \emph{De civ.\ Dei} IX.21).

Gregory Nazianzen joins the angels' knowledge to their joy. The Good
Shepherd lit the lamp of his flesh, swept the house, and found the lost
coin, ``And He calls together His Angel friends on the finding of the coin,
and makes them sharers in His joy, whom He had made to share also the
secret of the Incarnation'' (\emph{Or.}\ 38.14). The friends who knew the
secret rejoice when the coin is found.

The Catechism carries the same teaching through the whole economy: ``From
the Incarnation to the Ascension, the life of the Word incarnate is
surrounded by the adoration and service of angels'' (CCC~333).

\angeldossiersettled{\ST{III q.~8 a.~4}; \ST{III q.~59 a.~6}; with \ST{I
q.~57 a.~5} and \ST{I q.~64 a.~1 ad~4}.}
 {Church teaching that Christ is the centre of the angelic world, that the
  angels were created through and for him, and that they serve his saving
  plan (CCC~331--333; Col 1:16; Heb 1:14). Common teaching that Christ,
  also in his human nature, is set above every order of angels (Eph
  1:20--22; 1~Pet 3:22; Irenaeus, Chrysostom). Thomist position that the
  angels belong to the mystical body of which Christ is Head, that the
  grace of his humanity reaches their spirits by the personal union, and
  that his judicial power extends to them by reason of his human nature
  (III q.~8 a.~4; q.~59 a.~6). Thomist position that the angels knew the
  Incarnation in general from the beginning of their beatitude and learn
  its particulars by revelation (I q.~57 a.~5 ad~1).}
 {Col 2:10 (s.c.\ of III q.~8 a.~4); 1~Cor 6:3 (s.c.\ of III q.~59 a.~6);
  Eph 1:20--22; Heb 1:4--14, 2:5--16; Phil 2:10; Apoc 7:11; Matt 4:11;
  Dionysius, \emph{CH}~7.3 (Parker); Augustine, \emph{De vera rel.}\ 31 and
  \emph{De civ.\ Dei} IX.21, as cited by Aquinas; Irenaeus, \emph{Adv.\
  haer.}\ III.16.6; Chrysostom, \emph{Hom.\ in Heb.}\ 3.1, 5.1, \emph{Hom.\
  in Eph.}\ 3; Gregory Nazianzen, \emph{Or.}\ 38.14; CCC~331--333.}

\sectionguard
\subsection{The Annunciation (\ST{III q.~30})}\label{sec:cm-annunciation}

The Gospel names the messenger before it names the Virgin: ``the angel
Gabriel was sent from God into a city of Galilee, called Nazareth, To a
virgin espoused to a man whose name was Joseph, of the house of David: and
the virgin's name was Mary. And the angel being come in, said unto her:
Hail, full of grace, the Lord is with thee: blessed art thou among women''
(Luke 1:26--28). The same Gabriel had stood at the altar of incense six
months before and said to Zachary, ``I am Gabriel, who stand before God and
am sent to speak to thee and to bring thee these good tidings'' (1:19). To
Mary he announces the Son of the Most High, answers her question with the
overshadowing of the Holy Ghost, gives her the sign of Elizabeth, and
receives her answer: ``Behold the handmaid of the Lord: be it done to me
according to thy word. And the angel departed from her'' (1:38).

Dionysius reads the Annunciation as the law of the hierarchies at work in
the Incarnation itself: ``Angels first were initiated in the Divine mystery
of the love of Jesus towards man, then, through them, the gift of its
knowledge passed to us. Thus, for example, the most divine Gabriel
instructed Zachariah~\dots\ and he revealed to Mary, how, in her, should be
born the supremely Divine mystery of the unutterable God-formation''
(\emph{CH}~4.4; see \S\ref{sec:ch4}). Aquinas gives the mystery a question of
four articles.

\paragraph{Why it was announced (a.~1).} The objections hold the
announcement needless: the virginal conception was foretold by a prophecy of
predestination, which is fulfilled without our consent; the Virgin already
believed in the Incarnation; and those who conceive Christ spiritually
receive no such message. The \emph{sed contra} is the angel's ``Behold, thou
shalt conceive in thy womb'' (Luke 1:31). Aquinas gives four reasons. The
first is order: Mary should be instructed in mind before she conceived in
the flesh, for, as Augustine says, ``Mary is more blessed in receiving the
faith of Christ, than in conceiving the flesh of Christ.'' The second is
witness: instructed by God, she became a more certain witness of the
mystery. The third is the free gift of obedience, ``which she proved herself
right ready to do, saying: `Behold the handmaid of the Lord.'\,'' The
fourth is the marriage of God with mankind: ``there is a certain spiritual
wedlock between the Son of God and human nature. Wherefore in the
Annunciation the Virgin's consent was besought in lieu of that of the entire
human nature'' (\latin{loco totius humanae naturae}). The prophecy of
predestination ``is fulfilled without the causality of our will; not without
its consent'' (ad~1); the Virgin believed, but ``being humble, she did not
think such high things of herself'' (ad~2).

Leo the Great sees the same care of God for the Virgin's mind: ``lest in
ignorance of the heavenly counsel she should tremble at so strange a result,
she learns from converse with the angel that what is to be wrought in her
is of the Holy Ghost'' (\emph{Serm.}\ 21.1). Bernard, in the homilies
\emph{Super Missus est} that he wrote in praise of the Virgin Mother, turns
the angel's waiting into the waiting of the world:

\begin{quote}
The angel awaits your reply, for it is time that he should return to God,
Who sent him. We, too, are waiting, O Lady, for a word of mercy---we, who
are groaning under the sentence of condemnation. See, the price of our
salvation is offered to you; if you consent, we shall at once be delivered.
\dots\ Answer, then, quickly to the angel---yes, through the angel give
your consent to your God. Answer the word, receive the Word. Utter yours,
conceive the Divine. (\emph{Super Missus est} hom.~4; tr.\ 1909, pp.\
68--69)
\end{quote}

\noindent Adam, Abraham, and David, Bernard says, entreat her; ``the entire
human race'' lies at her feet in expectation. The angel stands between God
and the Virgin as the messenger who carries the question and will carry the
answer: ``through the angel give your consent to your God.''

\paragraph{Why by an angel (a.~2).} The objections: revelations to the
highest angels come from God immediately, and the Mother of God is above
the angels; divine things are taught to a woman by her husband; the highest
angels did not fully know the mystery, for they ask, ``Who is this that
cometh from Edom?''; and the greatest message should have come from the
highest order, whereas Gabriel is an archangel. The \emph{sed contra} is
Luke 1:26. Aquinas gives three reasons of fittingness. First, ``that in this
also might be maintained the order established by God, by which Divine
things are brought to men by means of the angels,'' and he quotes the
Dionysian text just given. Second, it befitted the restoration of mankind,
and he quotes Bede's homily for the Annunciation: the ruin began when the
devil sent the serpent to the woman, and the restoration begins when God
sends an angel to the Virgin; Eve, deceived through a serpent, offered her
husband the taste of death, and Mary, taught through an angel, brought forth
the author of salvation (Bede, \emph{Hom.}\ I.1; PL~94;
\S\ref{sec:fa-bede}). Third, it befitted her virginity, ``because
virginity is ever akin to the angelic nature. Surely to live in the flesh
and not according to the flesh is not an earthly but a heavenly life,''
words Aquinas takes from a sermon on the Assumption that circulated under
Jerome's name. Bede says the same in his own homily: she who strove to
imitate the angelic life deserved to enjoy both the sight and the speech of
an angel (\emph{Hom.}\ I.1).

The replies treat the Virgin's rank. ``The Mother of God was above the
angels as regards the dignity to which she was chosen by God. But as
regards the present state of life, she was beneath the angels'' (ad~1;
\latin{superior erat Angelis quantum ad dignitatem~\dots\ quantum ad statum
praesentis vitae, inferior erat Angelis}). Christ himself, by his passible
life, ``was made a little lower than the angels''; but he was wayfarer and
comprehensor at once and needed no angel's instruction, whereas his Mother
``was not yet in the state of comprehension.'' It was fitting that she learn
the mystery from an angel and not from a man, since she stood outside the
common rule of women (ad~2), and before Joseph, who learned it after she had
conceived. The angels knew the mystery, and their question from Edom is the
desire to learn its details, which no created intellect comprehends (ad~3).

On Gabriel's rank the doctors differ. ``Some say that Gabriel was of the
highest order,'' because Gregory writes that it was right for one of the
highest angels to come, since his message was the most sublime; Aquinas
holds that he was an archangel, \latin{de ordine Archangelorum}, as the
Church names him, and that ``it is therefore sufficiently credible that he was the
highest of the archangels''; his name, as Gregory says, means ``Power of
God,'' and fittingly, for ``the Lord of hosts and mighty in battle was
coming to overcome the powers of the air'' (ad~4). Bernard reads the same
name another way. ``I do not think that this was one of the lower angels
who for one cause or another are often sent to earth; and I gather it from
his name, which is interpreted `Strength of God'; because, also, he was not
sent as is usual from a superior spirit, but from God Himself'' (\emph{Super
Missus est} hom.~1). ``Among the blessed spirits themselves Gabriel alone
was excepted, for he alone was found worthy of his name and embassy.''
Bernard adds a pious conjecture of his own: the angel who strengthened
Joseph, ``a humble and timorous man,'' was very likely the same archangel
(\emph{ibid.}; Matt 1:20).

\paragraph{Why in bodily vision (a.~3).} The objections prefer an
intellectual vision, as more excellent, or an imaginary one, as Joseph's
was in sleep, and fear the troubling of the Virgin's mind. The \emph{sed
contra} is a sermon Aquinas reads as Augustine's, in which the Virgin says,
``To me came the archangel Gabriel with glowing countenance, gleaming robe,
and wondrous step.'' Aquinas determines that the angel appeared to her in a
bodily vision, for three reasons. The first concerns what was announced:
``the angel came to announce Incarnation of the invisible God. Wherefore it
was becoming that, in order to make this known, an invisible creature should
assume a form in which to appear visibly: forasmuch as all the apparitions
of the Old Testament are ordered to that apparition in which the Son of God
appeared in the flesh.'' The second concerns her dignity, since she was to
receive the Son of God ``not only in her mind, but in her bodily womb,'' and
so her bodily senses too were to be ``refreshed by the angelic vision''
(\latin{visione angelica refovendi}). The third is certainty, for which
Aquinas cites Chrysostom. Chrysostom's own question, on Joseph's dream, is
why the angel did not come ``openly, as to the shepherds, and to Zacharias,
and to the Virgin'': Joseph ``was exceedingly full of faith, and needed not
this vision. Whereas the Virgin, as having declared to her very exceeding
good tidings, greater than to Zacharias, and this before the event, needed
also a marvellous vision'' (\emph{Hom.\ in Matt.}\ 4.10). The replies add
that the Virgin had intellectual illumination together with the bodily
vision, though as a wayfarer she could not see the angel in his essence
(ad~1), and that the disturbance of a creature raised above itself does it
no harm; the angel at once says ``Fear not'' (ad~3). Aquinas records
another view as well: that the Virgin, accustomed to angelic visions, was
troubled not at seeing the angel but ``at his saying'' (ad~3; Luke 1:29).

Ambrose, whom the reply cites, draws from the scene a school of virginal
modesty: \latin{Trepidare virginum est, et ad omnes viri ingressus pavere,
omnes viri affatus vereri} (\emph{In Luc.}\ II.8)---it belongs to virgins to
tremble, to fear at every man's coming in, to be wary of every man's
address. He pictures the place: \latin{Sola in penetralibus, quam nemo
virorum videret, solus angelus reperiret: sola sine comite, sola sine teste;
ne quo degeneri depravaretur affatu, ab angelo salutatur}---alone in her
inner room, where no man might see her and only an angel could find her,
without companion and without witness, she is greeted by an angel, so that
no unworthy address might touch her. \latin{Disce, virgo, verborum vitare
lasciviam: Maria etiam salutationem angeli verebatur}: learn, virgin, to
shun wanton words; Mary was wary even of an angel's greeting. Bernard
enters the same room and finds it closed:

\begin{quote}
Closed, therefore, at that hour was the dwelling of that most prudent
Virgin, but to men, not to angels. For the angels are wont to be near those
who pray; they delight in beholding them raise their pure hands to heaven;
and with glad service they present to God the sacrifices of devotion which
they offer in the odour of sweetness. How pleasing to the Most High were
the prayers of Mary is well shown by the reverence with which the angel
saluted her. (\emph{Super Missus est} hom.~3; tr.\ 1909, p.~49)
\end{quote}

\noindent Bernard adds that no bolt could hold out ``the subtlety of his
nature,'' and that God, ``swifter than the angel,'' was already with the
Virgin when his messenger arrived: ``He Who had sent the angel was already
found by him with the Virgin.''

\paragraph{The order of the message (a.~4).} The objections find the order
unfitting: the angel praised her before announcing the Child, and gave
Elizabeth's conception, a lesser wonder, as proof of a greater. The
\emph{sed contra} is ``Those that are of God, are well ordered'' (Rom 13:1).
The angel had a threefold purpose: to draw her attention ``by a new and
unwonted salutation,'' in which ``Full of grace'' asserts her worthiness,
``The Lord is with thee'' announces the conception, and ``Blessed art thou
among women'' foretells her honour; to instruct her in the mystery by
foretelling the conception and birth, declaring the Child's dignity, and
making known the mode, ``The Holy Ghost shall come upon thee''; and to lead
her mind to consent by the example of Elizabeth and the argument from God's
omnipotence. ``To a humble mind nothing is more astonishing than to hear its
own excellence'' (ad~1). On whether Mary doubted, the Fathers seem to
differ. Ambrose says explicitly that she did not: ``She does not
doubt fulfilment when she asks how it shall be done''; a text Aquinas reads
as Augustine's says that the angel declared the possibility ``to Mary, in
doubt about the conception''; Aquinas reconciles them by calling it a doubt ``of wonder rather than of unbelief'' (ad~2). Ambrose, in his own
exposition, grounds the angelic embassy in the greatness of the mandate:
\latin{Tanti namque mandati mysterium non hominis fuit, sed angeli ore
promendum} (\emph{In Luc.}\ II.15)---the mystery of so great a charge was to
be spoken not by a man's mouth but by an angel's; and the angel announced an
old and barren woman's conception to the Virgin \latin{ut possibile Deo
omne quod ei placuerit, assereret} (II.19), to show that all that pleases
God is possible to him.

\paragraph{The Annunciation in the Church's praise.} The Byzantine Church
lets the archangel himself wonder. Her office for the feast sings that
Gabriel ``was sent from heaven to announce the glad tidings of conception
unto a Virgin; and when he was come to Nazareth, he mused within himself,
being astonished at the marvel: O, how can he who on high is beyond
comprehension be born of a Virgin!~\dots\ He upon whom the Six-winged and
the Many-eyed are unable to gaze deigneth, by his word alone, to become
incarnate of her'' (Hapgood, \emph{Service Book}, 1906, p.~201). The
troparion of the day joins the faithful to the angel's greeting: ``Gabriel
announceth the glad tidings of grace. Wherefore let us also cry aloud with
him unto the Birth-giver of God: Hail, thou that art full of grace, the Lord
is with thee!'' (p.~202). The Latin Church repeats the angel's message
morning, noon, and night in the Angelus, whose versicle begins ``The Angel
of the Lord declared unto Mary,'' and whose collect asks ``that as we have
known the Incarnation of Christ Thy Son by the message of an Angel, so, by
His Passion and Cross, we may be brought to the glory of His resurrection''
(\emph{Manual of Prayers}, Baltimore, 1889, pp.~55--56). The Second Vatican
Council names the scene in its chapter on the Mother of God: the Virgin of
Nazareth \latin{ab Angelo nuntiante, Dei mandato, ut ``gratia plena''
salutatur} (\emph{Lumen gentium} 56), greeted by the announcing angel, at
God's command, as full of grace.

\angeldossier{\ST{III q.~30 aa.~1--4}.}
 {Scripture teaches that the angel Gabriel was sent from God to the Virgin,
  announced the conception of the Son of God by the Holy Ghost, and received
  her consent (Luke 1:26--38). Church teaching that the Virgin was greeted
  by the angel at God's command and freely consented (\emph{Lumen gentium}
  56). Common teaching that the Annunciation was fittingly made by an angel,
  according to the order by which divine things come to men through the
  angels, and fittingly to a virgin (a.~2; Dionysius, \emph{CH}~4.4; Bede).
  Thomist position on the four reasons for the announcement, the bodily
  vision, and the order of the message (aa.~1, 3, 4). Disputed: Gabriel's
  rank (a.~2 ad~4). Pious tradition: Bernard's identification of the angel
  of Joseph's dream with Gabriel.}
 {Luke 1:31 (s.c.\ a.~1); Luke 1:26 (s.c.\ a.~2); a sermon read as
  Augustine's (s.c.\ a.~3); Rom 13:1 (s.c.\ a.~4); Dionysius, \emph{CH}~4.4; Augustine, \emph{De sancta virg.}\ 3, as cited; Bede, \emph{Hom.}\
  I.1 (PL~94); Chrysostom, \emph{Hom.\ in Matt.}\ 4.10; Ambrose, \emph{In
  Luc.}\ II.8, 15, 19; Gregory, \emph{Hom.\ in Ev.}\ 34, as cited; Leo,
  \emph{Serm.}\ 21.1; Bernard, \emph{Super Missus est} homm.~1, 3, 4;
  Byzantine office of the Annunciation (Hapgood 1906); the Angelus.}
 {On Gabriel's order: ``some'' place him in the highest order, from
  Gregory's phrase (a.~2 ad~4), and Bernard holds that he was sent from God
  himself and not through a superior spirit; Aquinas, the highest of the
  archangels. On the Virgin's question: Ambrose, no doubt; the text read as
  Augustine's, a doubt; Aquinas, a doubt of wonder (a.~4 ad~2). On her
  trouble: ``others'' say she was troubled at the words and not the sight
  (a.~3 ad~3).}

\sectionguard
\subsection{The Nativity, the Desert, and the Agony}\label{sec:cm-life}

\paragraph{The heavenly army at Bethlehem.} ``And behold an angel of the
Lord stood by them and the brightness of God shone round about them: and
they feared with a great fear. And the angel said to them: Fear not; for,
behold, I bring you good tidings of great joy that shall be to all the
people: For, this day is born to you a Saviour, who is Christ the Lord, in
the city of David.~\dots\ And suddenly there was with the angel a multitude
of the heavenly army, praising God and saying: Glory to God in the highest:
and on earth peace to men of good will'' (Luke 2:9--14). The Catechism
says that their song ``has not ceased resounding in the Church's praise''
(CCC~333).

Leo the Great preached on this night year after year, and he reads the
angels' joy as the measure of ours:

\begin{quote}
Therefore the exulting angel's song when the Lord was born is this, ``Glory
to God in the Highest,'' and their message, ``peace on earth to men of good
will.'' For they see that the heavenly Jerusalem is being built up out of
all the nations of the world: and over that indescribable work of the
Divine love how ought the humbleness of men to rejoice, when the joy of the
lofty angels is so great? (\emph{Serm.}\ 21.2)
\end{quote}

\noindent The angels rejoice because they see their own city being built out
of the nations. In another Christmas sermon Leo says that ``Today the
shepherds learnt from angels' voices that the Saviour was born in the
substance of our flesh and soul,'' so that ``we too may say with the army of
the heavenly host: `Glory in the highest to God, and on earth peace to men
of good will'\,'' (\emph{Serm.}\ 26.1), and he names peace as the gift
``which at the Lord's nativity was first proclaimed by the angel-choir''
(26.3).

Gregory Nazianzen ends his oration on the Theophany by calling his people
into the choir, ``With Shepherds glorify Him; with Angels join in chorus,''
for the heavenly hosts keep the feast with men ``because they love men, and
they love God'' (\emph{Or.}\ 38.17; \S\ref{sec:fb-nazianzen}): the feast is
one feast of heaven and earth. Earlier in the same oration Gregory
calls the Christ-child by Isaias's title in the Greek Bible, ``The Angel of
the Great Counsel of the Father'' (\emph{Or.}\ 38.2; Isa 9:6, where the
Septuagint reads ``the Messenger of great counsel,'' tr.\ Brenton): he who is the Lord of
the angels comes as the messenger of the Father's counsel.

Ambrose counts the embassies of that season as the care of God
establishing faith: \latin{angelus Mariam, angelus Ioseph, angelus pastores
edocet. Non satis est semel missum; duobus enim vel tribus testibus stat
omne verbum} (\emph{In Luc.}\ II.51)---an angel teaches Mary, an angel
Joseph, an angel the shepherds; one sending was not enough, for every word
stands on two or three witnesses. The heavenly host is rightly called an
army, \latin{qui ducem militiae sequebantur} (II.52), for they were
following the captain of their host. And the shepherds themselves are a
figure of the Church's guardians: \latin{non solum episcopos ad tuendum
gregem Dominus ordinavit, sed etiam angelos destinavit} (II.50)---the Lord
has appointed not only bishops to guard the flock but angels as well.
\latin{Vides pastores angelo credidisse} (II.53): the shepherds believed an
angel, and Ambrose bids his hearers believe the Father, the Son, the Holy
Spirit, the angels, the prophets, and the apostles.

The angel who appeared to Joseph ``in his sleep'' (Matt 1:20) warned him
again to fly into Egypt and called him back (Matt 2:13, 19). Chrysostom
explains the difference of the visions: the shepherds and the Virgin saw the
angel openly, while Joseph, ``exceedingly full of faith,'' needed no more
than a dream (\emph{Hom.\ in Matt.}\ 4.10). The Catechism sums up these
embassies of the Infancy: the angels ``protect Jesus in his infancy, serve
him in the desert, strengthen him in his agony in the garden'' (CCC~333).

\paragraph{Ministers in the desert.} ``Then the devil left him; and behold
angels came and ministered to him'' (Matt 4:11); Mark adds that ``he was
with beasts: and the angels ministered to him'' (Mark 1:13). Chrysostom
asks why they came only at the end:

\begin{quote}
For when the assault was going on, He suffered them not to appear, that He
might not thereby drive away the prey; but after He had convicted him in
all points, and caused him to take to flight, then they appear: that thou
also mayest learn, that after thy victories which are copied from His,
angels will receive thee also, applauding thee, and waiting as guards on
thee in all things. (\emph{Hom.\ in Matt.}\ 13.5)
\end{quote}

\noindent The angels who came to Christ after the victory are the angels
who will receive the Christian after his, as they carried Lazarus. Ambrose
sees the two natures in Mark's two clauses: \latin{Tuum est quod bestias
patitur, suum quod ab angelis praedicatur} (\emph{In Luc.}\ II.52)---that he
endures the beasts is yours, that he is proclaimed by angels is his own. For
Aquinas the same verse is a sign of the headship (\ST{III q.~8 a.~4}) and of
the dispensing that belongs to the Man Christ (\ST{III q.~59 a.~6}).
Gregory Nazianzen puts it among the feasts of Christ's life: ``you shall see
Him tempted and conquering and served by Angels'' (\emph{Or.}\ 38.16).

\paragraph{The angel of the agony.} ``And there appeared to him an angel
from heaven, strengthening him. And being in an agony, he prayed the
longer'' (Luke 22:43). At his arrest Christ tells Peter that he could ask
his Father, ``and he will give me presently more than twelve legions of
angels'' (Matt 26:53). Hilary of Poitiers treats the angel of the agony in
his work on the Trinity. He reports that ``in many manuscripts, both Latin
and Greek, nothing is said of the angel's coming or the Bloody Sweat,'' and
leaves the question of the text in suspense; but he forbids the heretics to
find in the verse a proof of weakness:

\begin{quote}
Let them remember the Creator of the angels needs not the support of His
creatures. Moreover His comforting must be explained in the same way as His
sorrow. He was sorrowful for us, that is, on our account; He must also have
been comforted for us, that is, on our account. If He sorrowed concerning
us, He was comforted concerning us. (\emph{De Trin.}\ X.41)
\end{quote}

\noindent Hilary reads the angel as sent ``to watch over the Apostles,'' so
that the Lord, comforted concerning them, could say, ``Sleep on now, and
take your rest'' (X.40). Chrysostom gives the verse the opposite weight: the
sweat like drops of blood and the strengthening angel were ``a thousand sure
signs of fear, lest any one should affirm the words to be feigned''
(\emph{Hom.\ in Matt.}\ 83.1). Hilary guards the Godhead, Chrysostom the
reality of the manhood; both read the angel as serving the salvation of the
disciples and of the Church. Ambrose, at the arrest, sets the legions aside
before the Lord's word: \latin{Quo mihi legiones angelorum, quo coelestis
exercitus? Vox Domini sola plus terret} (\emph{In Luc.}\ X.65)---what need
of legions of angels or a heavenly army? The Lord's voice alone struck
greater fear, when those who came to seize him fell to the ground (John
18:6).

\sectionguard
\subsection{The Resurrection, the Ascension, and the Coming in Glory}\label{sec:cm-glory}

\paragraph{The witnesses of the tomb.} ``For an angel of the Lord descended
from heaven and coming rolled back the stone and sat upon it. And his
countenance was as lightning and his raiment as snow.~\dots\ And the angel
answering, said to the women: Fear not you: for I know that you seek Jesus
who was crucified. He is not here. For he is risen, as he said'' (Matt
28:2--6). Luke has two men ``in shining apparel'' who ask, ``Why seek you
the living with the dead?'' (Luke 24:4--5), and John has Magdalen see ``two
angels in white, sitting, one at the head, and one at the feet, where the
body of Jesus had been laid'' (John 20:12). Cyril of Jerusalem numbers the
angels among the witnesses of the Resurrection: ``Angels of God who were
present testified of the Resurrection of the Only-begotten,'' and the women
``beheld the mighty earthquake, and the radiance of the Angel who stood
by'' (\emph{Cat.}\ 14.22). Ambrose finds in the differing numbers of the
Gospels a sign of the angels' service: \latin{Diversitas autem visionum,
frequentiam famulantium significat angelorum} (\emph{In Luc.}\
X.181)---the diversity of the visions signifies the frequency of the angels
who serve, as the Lord promised that the disciples would see ``the angels of
God ascending and descending upon the Son of man'' (John 1:51). Gregory
Nazianzen bids his hearers go to the tomb with the women: ``Be first to see
the stone taken away, and perhaps you will see the Angels and Jesus
Himself'' (\emph{Or.}\ 45.24).

\paragraph{The escort of the Ascension.} ``And while they were beholding
him going up to heaven, behold two men stood by them in white garments. Who
also said: Ye men of Galilee, why stand you looking up to heaven? This Jesus
who is taken up from you into heaven, shall so come as you have seen him
going into heaven'' (Acts 1:10--11). Leo draws from the two angels the
whole service of the angels to the Incarnate Word, from the announcing of
his conception to the foretelling of his coming to judge, ``that we might
understand what great powers will come with Him as Judge, when such great
ones ministered to Him even in being judged'' (\emph{Serm.}\ 74.4); in the
Ascension the nature of mankind passed above the angels' ranks and the
archangels' heights (\emph{Serm.}\ 73.4;
\S\ref{sec:fl-leo-gennadius}). The feast commemorates the day ``on which the Nature of our humility in
Christ was raised above all the host of heaven, over all the ranks of
angels, beyond the height of all powers, to sit with God the Father''
(\emph{Serm.}\ 74.1).

The Fathers read the entrance psalm, ``Lift up your gates, O ye princes, and
be ye lifted up, O eternal gates: and the King of Glory shall enter in. Who
is this King of Glory?'' (Ps 23[24]:7--8), as the dialogue of the angels at
the Ascension. Cyril reminds his candidates ``that the divine powers also
said to one another, Lift up your gates, ye Princes'' (\emph{Cat.}\ 14.24).
Gregory Nazianzen makes his hearers part of the escort:

\begin{quote}
And if He ascend up into Heaven, ascend with Him. Be one of those angels who
escort Him, or one of those who receive Him. Bid the gates be lifted up, or
be made higher, that they may receive Him, exalted after His Passion. Answer
to those who are in doubt because He bears up with Him His body and the
tokens of His Passion, which He had not when He came down, and who therefore
inquire, ``Who is this King of Glory?'' that it is the Lord strong and
mighty. (\emph{Or.}\ 45.25)
\end{quote}

\noindent Gregory joins the psalm to the question of Isaias's drama, ``Who
is this that cometh from Edom and from the things of earth? Or How are the
garments red of Him that is without blood or body, as of one that treads in
the full wine-press?'' and answers with ``the beauty of the array of the
Body that suffered, adorned by the Passion, and made splendid by the
Godhead'' (\emph{ibid.}; Isa 63:1--2). The angels who asked at the gates are
the angels whom Dionysius shows asking Jesus the same question, and whom
Jesus teaches (\emph{CH}~7.3). The doubt of the gatekeepers is wonder at the
wounds.

\paragraph{The angels of the Judgment.} ``And when the Son of man shall come
in his majesty, and all the angels with him, then shall he sit upon the seat
of his majesty'' (Matt 25:31). ``He shall send his angels with a trumpet and
a great voice: and they shall gather together his elect from the four
winds'' (Matt 24:31); ``The Son of man shall send his angels, and they shall
gather out of his kingdom all scandals, and them that work iniquity'' (Matt
13:41); ``the Lord himself shall come down from heaven with commandment and
with the voice of an archangel and with the trumpet of God'' (1~Thess
4:15[16]). Cyril sets the two comings side by side: ``In His first coming,
He endured the Cross, despising shame; in His second, He comes attended by a
host of Angels, receiving glory,'' so that ``when with Angels we meet our
Master, we may worship Him and say, Blessed is He that cometh in the Name of
the Lord'' (\emph{Cat.}\ 15.1). He describes the Judge ``escorted by tens of
thousands of Angels'' (15.10), the sign of the Cross going before, the
enemies of Christ unable to flee because ``the Angel hosts shall encompass
them,'' and the elect gathered by the angels: ``That King, so great and
glorious, attended by the Angel-guards, the partner of the Father's throne,
will not despise His own servants'' (15.22). The Archangel will ``make
proclamation and say to all, Arise to meet the Lord'' (15.21). Aquinas,
answering the objection that the angels who come to judge cannot be judged,
keeps both truths: the angels judge as spiritual creatures, and are judged
by Christ, who is the Truth (\ST{III q.~59 a.~6 ad~2}).

\sectionguard
\subsection{Mary, Queen of Angels}\label{sec:cm-queen}

\paragraph{Above the choirs.} The angel who saluted the Virgin was greeting
his Queen. Damascene says it in a sentence: ``Beautiful is the angel's
salutation to her who is greater than an angel'' (\emph{Hom.\ in Dorm.}\
I). Bernard, in his first homily on the Annunciation, reasons from her Son
to her rank:

\begin{quote}
Do you not judge, and rightly, that she who has the God-man for her Son is
exalted in greatness above all the choirs of angels? Did not Mary
confidently call the God and Lord of Angels her Son, saying: ``Son, why
hast thou done so to us?'' Which of the angels would have presumed thus to
speak? It is sufficient for them and something great, that while by nature
they are spirits by grace they are made and called angels.~\dots\ In
confidently calling God her Son, Mary acknowledges herself mother of that
Majesty Whom those angels serve with reverential awe. (\emph{Super Missus
est} hom.~1; tr.\ 1909, p.~30)
\end{quote}

\noindent ``God to Whom the angels are subject. God, Whom the powers and
principalities obey, was subject to Mary'' (p.~31; Luke 2:51). In his second
homily Bernard calls her the one ``foreknown by the Most High, prepared for
Himself, guarded by angels, pointed out by the Patriarchs, promised by the
Prophets.'' Aquinas states the same exaltation with its limit: she was above
the angels ``as regards the dignity to which she was chosen by God,'' and
beneath them in the state of this life, since she was not yet a
comprehensor (\ST{III q.~30 a.~2 ad~1}).

The Church teaches this exaltation in her acts. Pius IX's bull
\emph{Ineffabilis Deus}, as Pius XII quotes it, says that God endowed
her with heavenly gifts \latin{longe ante omnes angelicos spiritus cunctosque
Sanctos}---far beyond all the angelic spirits and all the saints---and calls
her Queen of heaven and earth, \latin{super omnes Angelorum choros
Sanctorumque Caelitum ordines exaltata}, exalted above all the choirs of
angels and the orders of the saints (\emph{Ad caeli Reginam} 41--42). The
Second Vatican Council says that Mary, \latin{per gratiam Dei post Filium
prae omnibus angelis et hominibus exaltata}, exalted by the grace of God
after her Son above all angels and men, is rightly honoured with a special
cult (\emph{Lumen gentium} 66), and bids the faithful pray to her, \latin{nunc
quoque in coelo super omnes Beatos et Angelos exaltata}, now exalted in
heaven above all the blessed and the angels (69). Pius XII's encyclical on
her queenship, \emph{Ad caeli Reginam} (11 October 1954), opens with the
Assumption, by which she reigns in soul and body \latin{inter Angelorum
Sanctorumque Caelitum choros} with her only-begotten Son (3), and closes
with the Roman Breviary's words for the Assumption, that she is
\latin{exaltata super choros angelorum ad caelestia regna}, exalted above the
choirs of angels to the heavenly kingdom (46). By the same letter Pius XII
instituted the feast of Mary's Queenship for 31 May (47).

\paragraph{Queen of Angels in the Church's prayer.} The title stands in the
Litany of Loreto after the Help of Christians and before the Queen of
Patriarchs: \latin{Regina Angelorum}, ``Queen of Angels, pray for us''
(\emph{Manual of Prayers}, 1889, p.~67). Pius XII names the Litany among the
witnesses that daily invite the Christian people to call upon Mary as Queen
(\emph{Ad caeli Reginam} 31). Among the Marian antiphons that the Latin Church sings (\emph{Ad caeli
Reginam} 30), the \latin{Ave Regina caelorum} greets her by the same title:
\latin{Ave, Regina caelorum! Ave, domina angelorum!} The Baltimore
\emph{Manual} appoints it from the Purification to Maundy Thursday and
renders it, ``Hail! O Queen of Heav'n enthron'd! Hail! by angels
mistress own'd!'' (p.~77). The Missal of 1962 opens the Mass of the Queen,
on 31 May, with the Introit \latin{Gaudeamus omnes in Domino, diem festum
celebrantes sub honore beatae Mariae Virginis Reginae: de cuius solemnitate
gaudent Angeli, et collaudant Filium Dei}: let us all rejoice in the Lord,
keeping the feast in honour of the Blessed Virgin Mary, Queen, at whose
solemnity the angels rejoice and praise together the Son of God
(\emph{Missale Romanum}, 1962, 31 May, n.~2681).

The Byzantine Church sings the same truth at Vespers and at the dismissal of
the Divine Liturgy in her hymn to the Mother of God: ``More honourable than the
Cherubim, and beyond compare more glorious than the Seraphim, thou who
without defilement barest God the Word, true Birth-giver of God, we magnify
thee'' (Hapgood, \emph{Service Book}, pp.~14, 126). The two highest orders of
Dionysius's first hierarchy are named in order to set her above them.
Damascene turns the same comparison into images drawn from the Old
Testament: ``Thou art the royal throne which angels surround, seeing upon it
their very King and Lord''; and she is Jacob's ladder, on which angels
descended to the Word made flesh, ``ministering to Him as their God and
Lord, and men, adopting the life of angels, are carried up to heaven''
(\emph{Hom.\ in Dorm.}\ I; Gen 28:12).

\paragraph{The angels at her Assumption.} Pius XII defined, on 1 November
1950, that the Immaculate Mother of God, \latin{expleto terrestris vitae
cursu, fuisse corpore et anima ad caelestem gloriam assumptam}---having
completed the course of her earthly life, was assumed body and soul into
heavenly glory (\emph{Munificentissimus Deus}). Among the reasons the
theologians gave, he records her ``exalted holiness, entirely surpassing
the sanctity of all men and of the angels,'' and he quotes Albert the
Great's conclusion that ``the most blessed Mother of God has been assumed
above the choirs of angels'' (\latin{super choros Angelorum est
assumpta}). The Fathers and the liturgy fill in the angels' part.
Damascene sets them about her death and her passage:

\begin{quote}
Angels with archangels bear thee up. Impure spirits trembled at thy
departure. The air raises a hymn of praise at thy passage, and the
atmosphere is purified. Heaven receives thy soul with joy. The heavenly
powers greet thee with sacred canticles and with joyous praise~\dots\
(\emph{Hom.\ in Dorm.}\ I; tr.\ Allies, 1898)
\end{quote}

\noindent In his second homily all nine orders keep the feast: ``Angels and
Archangels are in jubilation, Powers exult, Principalities and Dominations,
Virtues and Thrones are in gladness: Cherubim and Seraphim magnify God. Not
the least of their praise is it to refer praise to the Mother of glory''
(\emph{Hom.\ in Dorm.}\ II). He hears ``the swelling hymns of angels who
precede, accompany, and follow her,'' some forming ``the guard of honour to
that undefiled and immaculate~\dots\ soul on its way to heaven until the queen
reaches the divine throne,'' others surrounding the body and proclaiming
the Mother of God ``in angelic harmony'' (\emph{ibid.}).

Theodore the Studite, in his encomium on the Dormition, describes the
Lord's coming for his Mother \gk{meta dox\=es ischyos autou kai pas\=es
stratias ouranou}, with the glory of his might and all the army of heaven,
and the mingled liturgy of the two worlds: \gk{kai aorat\=os
men eleitourgoun hoi as\=omatoi, s\=omatik\=os de hoi apostoloi} (\emph{Or.}\
5.5; PG~99, 728)---the bodiless ministered invisibly, the apostles in the
body; \gk{symmig\=es \=en, adelphoi, h\=e pan\=egyris kai ho thiasos
ouranios te kai epigeios}, mingled, brethren, was the festival and the
company, heavenly and earthly. He names them all, Angels, Archangels,
Dominions, Thrones, Principalities, Authorities, Powers, the Cherubic and
Seraphic orders, with apostles, martyrs, and the just, some running before
and some coming to meet her, some leading and some following, all crying
with one voice in gladness. The Byzantine office for the feast gives the
angels their words at the gates:

\begin{quote}
The heavenly powers most high also, when they were come with their Master,
were seized with dread as they escorted the all-pure body which had received
God; with stately mien, also, they went before, and invisibly cried aloud
unto the Powers most high: Lo, the Queen over all, and the Maiden of God
cometh. Be ye lifted up, O gates~\dots\ (Hapgood, \emph{Service Book}, p.~264)
\end{quote}

\noindent Its refrain sings, ``When the all-holy Angels beheld thine
Assumption they marvelled how a Virgin should ascend from earth into
heaven'' (p.~266). The Roman Missal sings the same joy in the Alleluia of
the feast: \latin{Assumpta est Maria in caelum: gaudet exercitus Angelorum}
(\emph{Missale Romanum}, 1962, 15 August, n.~3382)---Mary is taken up into
heaven; the army of the angels rejoices. The psalm of the gates that the
Fathers sang of Christ's Ascension is sung again of his Mother's.

\sectionguard
\subsection{The Redeemed among the Angelic Ranks}\label{sec:cm-ranks}

\paragraph{Equal to the angels.} The Lord promised that the children of the
resurrection ``are equal to the angels and are the children of God'' (Luke
20:36) and ``shall be as the angels of God in heaven'' (Matt 22:30). Aquinas
asks whether men are taken up into the angelic orders and answers with the
distinction that governs his whole account of the orders: ``Considered only
as regards the grade of nature, men can in no way be assumed into the
angelic orders; for the natural distinction will always remain''; but ``by
the gift of grace men can merit glory in such a degree as to be equal to the
angels, in each of the angelic grades; and this implies that men are taken
up into the orders of the angels'' (\ST{I q.~108 a.~8}; see \S\ref{sec:q108}).
He rejects the opinion that only virgins or the perfect are assumed, since
Augustine teaches that ``there will not be two societies of men and angels,
but only one; because the beatitude of all is to cleave to God alone''
(\emph{ibid.}; \emph{De civ.\ Dei} XII.9). Augustine himself teaches that
the restored part of mankind ``should fill up the gap which the rebellion
and fall of the devils had left in the company of the angels''
(\emph{Enchir.}\ 29), and that ``the holy angels, taught by God~\dots\ know
how great a number of the human race are to supplement their ranks'' (62;
see \S\ref{sec:fl-augustine}).

\paragraph{Gregory's tenth coin.} Gregory the Great gives the doctrine its
images in his homily on the lost sheep and the lost coin
(\S\ref{sec:fa-gregory}). The hundred sheep are the whole
rational creation, angels and men; the ninety-nine left in the desert are
the highest choirs of angels, and the lost sheep was sought on earth
\latin{ut perfecta summa ovium integraretur in coelo}, so that the perfect
sum might be made whole in heaven (\emph{Hom.\ in Ev.}\ 34.3). The woman's
ten coins say the same in another figure: nine orders of angels, and man
created the tenth to complete the number of the elect (34.6). Cyril of Jerusalem uses the parable in the same sense at the
Judgment: however great the multitude of mankind, ``it is little, for the
Angels are many more. They are the ninety and nine sheep, but mankind is
the single one'' (\emph{Cat.}\ 15.24).

Gregory then turns the orders into a mirror of Christian lives: as many
men will ascend \latin{quantos illic contigit electos angelos remansisse},
and they are assigned to the orders \latin{per conversationis
similitudinem}, by likeness of life, from the faithful announcers of little
things among the Angels to those who burn with contemplation among the
Seraphim, and he bids his hearers search themselves for their lot
(34.11--12). Aquinas uses this homily for the meaning of the names of the orders
(\ST{I q.~108 a.~5}) and sets its order of the orders beside Dionysius's
(a.~6); its teaching on the lot of men among the orders is the teaching of
a.~8.

\paragraph{A new angel.} Gregory Nazianzen calls man, set between the
visible and the invisible creation, ``a new Angel, a mingled worshipper,
fully initiated into the visible creation, but only partially into the
intellectual'' (\emph{Or.}\ 38.11). Chrysostom promises the conqueror of
temptation that the angels who came to Christ in the desert ``will receive
thee also, applauding thee'' (\emph{Hom.\ in Matt.}\ 13.5). Damascene draws
from the Incarnation a boldness of his own. ``God did not unite Himself to
the angelic nature, but to the human. He did not become an angel,'' and our
nature, ``In itself~\dots\ far removed from the angels, on account of death
and the heaviness of the body,'' has ``through God's goodness and its union
with Him~\dots\ become higher than the angels. For angels stand by that
nature with fear and trembling, as, in the person of Christ, it sits upon a
throne of glory'' (\emph{De imag.}\ III; tr.\ Allies, 1898). He finds the
ground of this in the Eucharist, by which men partake of the Body united to
the Godhead: ``How are we not greater than the angels, if through fidelity
to the commandments we keep this perfect union?'' The Greek line after him
carries the comparison further, and Palamas, weighing man's image against
the angels', finds it more fully in man (see \S\ref{sec:bz-palamas}).

\paragraph{Already in their company.} The Church does not wait for the
resurrection to join the angels. ``You are come to mount Sion and to the
city of the living God, the heavenly Jerusalem, and to the company of many
thousands of angels'' (Heb 12:22). The Roman Mass ends each preface by
singing \latin{cum Angelis et Archangelis, cum Thronis et Dominationibus,
cumque omni militia caelestis exercitus}, with the angels and archangels,
with the thrones and dominations, and with the whole army of the heavenly
host (\emph{Missale Romanum}, 1962, Preface of Easter), and the Canon asks
that the offering be carried \latin{per manus sancti Angeli tui in sublime
altare tuum}, by the hands of thy holy angel to thine altar on high. The
Byzantine priest prays at the Little Entrance to the Lord ``who hast
appointed in heaven ranks and hosts of Angels and Archangels for the
ministry of thy glory: Cause that with our entrance may enter also the holy
Angels with us serving thee, and with us glorifying thy goodness''
(Hapgood, \emph{Service Book}, p.~83). The company that the redeemed will
complete is already present at their altars (see \S\ref{sec:byzantine}).
